\section{What literature disagreement can establish about a benchmark}
\label{exec17:validity-development}

A benchmark graph is a scientific object with a population, measurement scheme
and intended regime. Its usefulness does not follow from how often it has been
reused. Yet a new paper that appears to contradict an arrow is not automatically
a correction to that graph. In banking, an association between a credit offer
and repayment can concern applicants, approved borrowers or customers who
accepted the offer. A different selection rule changes the population law.
An account-level measurement can also differ from a customer-level mechanism.
Before an assistant votes for or against a diagram, it must establish what
proposition the diagram and each retrieved result actually assert.

Zhang and colleagues offer a useful starting point: retrieve domain papers,
derive association queries from a benchmark graph, and ask a language model
whether each paper supplies relevant empirical evidence. Their June 2026
preprint reports a survey of 96 papers using 69 benchmarks, followed by an
evaluation of 11 graphs using 38,081 full texts. The approach makes benchmark
assumptions objects of investigation rather than treating every inherited
diagram as unquestionable truth. Its statistical and graphical interfaces
nevertheless need repair before the proposed library could use the result to
change a causal model. The assessment here concerns the inspected first-version
text, pseudocode and prompt; it does not claim to reproduce the upstream
implementation.~\citep{exec17_benchmark_consistency}

\subsection{A graph implication is not its converse}

For a distribution $P$ that is Markov with respect to a DAG $G$, graphical
separation supplies an independence implication:
\begin{equation}
 X\perp_G Y\mid Z \quad\Longrightarrow\quad X\perp_P Y\mid Z.
\label{exec17:validity-markov}
\end{equation}
Here $\perp_G$ denotes d-separation and $\perp_P$ statistical conditional
independence. Reversing this implication requires a further assumption, such
as faithfulness: the distribution has no extra independences beyond those
implied by the graph. Even dependence need not be captured by a particular
linear correlation coefficient. The paper's description that unseparated
variables, including adjacent variables, are always correlated therefore goes
beyond the Markov premise.~\citep{exec17_benchmark_consistency}

An exact project example shows that positive density does not repair the
converse. Let independent disturbances $\varepsilon_X,\varepsilon_A,
\varepsilon_Y$ be standard Gaussian and define
\begin{equation}
 X=\varepsilon_X,\qquad A=X+\varepsilon_A,\qquad
 Y=X-A+\varepsilon_Y.
\label{exec17:validity-cancellation}
\end{equation}
The graph contains $X\to A$, $X\to Y$ and $A\to Y$, and the joint law is
Markov for this graph. In the order $(X,A,Y)$ its covariance matrix is
\begin{equation}
 \Sigma=\begin{pmatrix}1&1&0\\1&2&-1\\0&-1&2\end{pmatrix},
 \qquad \det(\Sigma)=1,\qquad
 \operatorname{Cov}(X,Y\mid A)=\tfrac12.
\label{exec17:validity-covariance}
\end{equation}
Indeed $Y=-\varepsilon_A+\varepsilon_Y$, so $X$ and $Y$ are marginally
independent despite their adjacency. The independent-disturbance coefficient
map is nonsingular, establishing positive density. The direct and mediated
effects cancel for this population; conditional association need not cancel.
The covariance and determinant are reproduced by exact coefficient arithmetic,
not by a sample whose correlations happen to be small.
\citep{exec17_advice_validity_exact}

For an offer, subsequent account action and final outcome, such a construction
distinguishes a mechanism from its total marginal association. It is not an
assertion that the bank has these coefficients. Its role is to show why a
literature result reporting independence cannot automatically invalidate an
adjacency under a Markov-only contract. A faithful-model comparison is a
different, explicitly stronger exercise. Conversely, an estimated nonzero
correlation does not by itself identify a direct arrow. A benchmark evaluator
must label which of these propositions it tests.

\subsection{Conditioning can open a path that a hitting set overlooks}

The published Algorithm 1 enumerates collider-free paths and returns minimal
sets that intersect each such path as conditional-independence sets. Blocking
paths that are initially active is necessary, but does not ensure that
conditioning leaves every other path blocked. A collider becomes active when
it, or a descendant, is conditioned upon. The paper's own d-separation
definition includes this condition; the inspected hitting-set pseudocode does
not apply it to candidate separators.~\citep{exec17_benchmark_consistency}

Consider the following five-node project DAG:
\begin{equation}
 X\to C,\quad B\to C,\quad C\to Z,\quad X\to Z,
 \quad Z\to Y,\quad B\to Y.
\label{exec17:validity-graph}
\end{equation}
Its collider-free paths from $X$ to $Y$ are $X-Z-Y$ and $X-C-Z-Y$. The unique
minimal hitting set of these paths is $\{Z\}$. Conditioning on $Z$ blocks
both, but activates $X\to C\leftarrow B\to Y$ through the descendant
$Z$ of collider $C$. It also activates the path
$X\to Z\leftarrow C\leftarrow B\to Y$ by conditioning on collider $Z$
itself. Thus $\{Z\}$ is not a d-separator. Enumerating every simple path
and applying the full activity criterion gives $\{B,Z\}$ as the minimal
d-separator. This is a counterexample to the published specification, not an
observed output of uninspected code.~\citep{exec17_advice_validity_exact}

A nonsingular Gaussian construction verifies the practical consequence:
\begin{equation}
\begin{split}
 X&=\varepsilon_X,\qquad B=\varepsilon_B,\qquad
 C=X+B+\varepsilon_C,\\
 Z&=X+C+\varepsilon_Z,\qquad Y=Z+B+\varepsilon_Y,
\end{split}
\label{exec17:validity-opened-scm}
\end{equation}
with independent standard Gaussian disturbances. Its coefficients give
$\operatorname{Var}(Z)=7$, $\operatorname{Cov}(X,Y)=2$,
$\operatorname{Cov}(X,Z)=2$ and $\operatorname{Cov}(Z,Y)=8$. Therefore
\begin{equation}
 \operatorname{Cov}(X,Y\mid Z)=2-\frac{2\times8}{7}=-\frac27.
\label{exec17:validity-opened-covariance}
\end{equation}
The exact covariance agrees with the independent active-path enumeration.
\citep{exec17_advice_validity_exact}

The remedy is concrete. Candidate conditioning sets must pass a sound
d-separation procedure for the stated graph family, including collider and
descendant activation. Small complete path enumeration is an independent
reference for testing a scalable implementation; it is not a scalable method
for every bank graph. If a graph has latent variables or selection nodes, the
evaluator needs the appropriate representation and separation criterion,
rather than treating its observed mixed graph as an ordinary DAG. The query
also needs to distinguish conditioning on a recorded covariate from
intervening on it. A controlled experiment and an observational subgroup do
not generally describe the same probability distribution.

\subsection{A nonsignificant test does not certify independence}

The inspected prompt instructs the verifier to infer independence when a
correlation test is not significant. That rule confuses a sample decision
with a population property. It can turn an underpowered study into apparent
evidence against a correct causal mechanism. It also treats a selected test
as if it were sensitive to every form of dependence.
\citep{exec17_benchmark_consistency}

For a finite illustration, let $X$ be a randomized binary treatment with
probability $1/2$ and let an independent flip $V$ have probability $1/4$.
Set $Y=X\mathbin{\mathrm{xor}}V$. The four population cell probabilities are
\begin{equation}
 (p_{00},p_{01},p_{10},p_{11})=\tfrac18(3,1,1,3),
 \qquad E{Y\mid\operatorname{do}(X=1)}
       -E{Y\mid\operatorname{do}(X=0)}=\tfrac12.
\label{exec17:validity-binary}
\end{equation}
In four independent observations, the event of one observation in every cell
has probability
\begin{equation}
 4!\left(\frac38\right)^2\left(\frac18\right)^2
       =\frac{27}{512}.
\label{exec17:validity-balanced}
\end{equation}
For that balanced table, Fisher's two-sided exact test has $p=1$. Conditional
on both margins being $(2,2)$, the upper-left cell under independence takes
values $0,1,2$ with probabilities $(1,4,1)/6$; all three tables have probability
no greater than that of the observed central table. The resulting sum is one.
Complete enumeration of all 256 four-observation sequences finds the 24
balanced sequences and the probability above. Yet the true randomized effect
remains $1/2$.~\citep{exec17_advice_validity_exact}

If an outcome contributes a stipulated USD 100 and treatment costs USD 30, the
population incremental value is $100(1/2)-30=\text{USD }20$ per decision.
Neither the four-row sample nor its large $p$ value establishes this value; it
is known here only because the example stipulates the structural law. The
example demonstrates the possible loss from converting a nonsignificant
finding into a zero-effect declaration. It is not a proposal ROI estimate.

A repaired extraction should preserve the estimator, sample size, sampling
design, uncertainty interval, null hypothesis and covariates. ``Did not reject
this null'' is an admissible extracted finding. ``Established independence''
requires a different justification. For a decision, an equivalence analysis
could target a scientifically chosen negligible-effect region, with its own
design and error control. A language model may help locate those details; it
cannot supply missing power, overlap or identification by rewriting the prose.

\subsection{Unknown answers change the target of an inconsistency rate}

The paper retains an unknown category and excludes unknown answers from its
inconsistency denominator. This is a useful distinction from forcing every
paper into a binary label, but it changes what a reported rate means. Write
$R=1$ when a query receives an answer and $D=1$ when that answer is categorized
as a contradiction. If $r=P(R=1)$ and $q=P(D=1\mid R=1)$, then without an
assumption on unknown cases, a latent full-query contradiction fraction $d$
obeys only
\begin{equation}
 rq\leq d\leq rq+1-r.
\label{exec17:validity-unknown}
\end{equation}
The bound assumes that a meaningful binary contradiction status exists for
every target query and that answered labels are correct; context mismatch or
misclassification requires further treatment. The original answered-case
metric does not make either additional premise true.
\citep{exec17_benchmark_consistency}

In an authored example with 20 contradictions among 30 answered queries out
of 100, the answered-case rate is $2/3$, the answer rate is $0.3$, and the
full-query range is $[0.2,0.9]$. Both endpoints are attained by assigning all
unknown cases to one category. This is a missing-information bound, not an
interval for statistical sampling error.~\citep{exec17_advice_validity_exact}
The pipeline should therefore report answer coverage alongside disagreement,
and distinguish unique papers, paper-pair queries and graph propositions.
Several queries can share one paper, and several papers can share a dataset;
their labels are not automatically independent replications.

Retrieval introduces another selection layer. The inspected procedure uses
PubMed/PMC and OpenAlex-related access routes, caps ranked search results at
500 per pair, and caps examined pairs at 200. Its verifier is the historical
Qwen3-30B-A3B-Thinking-2507 model, reported on an H100 80GB. Accessible text,
search ranking, variable synonyms and the relevance filter define the resulting
evidence pool. An inconsistency rate in that pool is not a random-sample
estimate of all scientific evidence or all graph errors. The authors report
90 correct judgments in a manual sample of 100; the linked detailed artifact
has not been inspected here, so this remains their reported validation.
\citep{exec17_benchmark_consistency}

There are also unresolved source-accounting questions. The main text lists
69 benchmarks while the appendix enumeration reaches 72. The introduction
describes Barley as misaligned, but the evaluated set excludes it. A separate
bibliometric search reports 11,633 LLM causal-discovery papers in 2025, while
its appendix includes MATH, ImageNet and COCO among prominent benchmark
labels. Without resolving search specificity, that count cannot establish
the size of the causal-discovery field. These limitations do not erase the
value of evidence retrieval; they limit the numerical and substantive claims
that this proposal can inherit.~\citep{exec17_benchmark_consistency}

\subsection{A repair that retains the useful human workflow}

The useful adoption target is a reviewed evidence queue. Each item should link
an exact passage to the original PDF and record population, unit, variable
definition, measurement date, conditioning set, intervention and selection
mechanism. A proposed conflict should identify both propositions and explain
why they refer to a comparable regime. For a banking diagram, a relationship
among approved borrowers cannot silently replace one among all applicants.
An updated sensor or accounting definition can require a new variable and
measurement model rather than an arrow reversal.

The proposal can conditionally adopt retrieval, synonym suggestions and
abstaining evidence extraction. It should reject automatic benchmark relabeling
from literature votes, adjacency-as-correlation rules and nonsignificance-as-
independence prompts. Direct software adoption or a port remains contingent on
inspecting pinned code, dependencies and rights. The mathematical repairs are
known: full separation checks, faithful versus Markov-only query types, correct
statistical statements, explicit missing-answer targets and reviewed context
alignment. Whether a repaired pipeline attains useful accuracy and cost is a
prototype question, not a result established by the present manuscript.

A public prototype can make that question testable before a bank trial. Build
a small adjudicated set of original passages and graph queries, including
known collider cases, cancellation, nonsignificant dependent samples, changed
populations and deliberately unavailable evidence. Keep both the graph truth
for constructed cases and the expert justification for real passages separate
from the model's suggestions. Compare keyword retrieval, human-only review and
an abstaining assistant on factual extraction, query validity, answer coverage,
false conflict alerts, review time and decision-relevant consequences. Preserve
disagreements instead of training on the test judgments. The mathematical
examples specify failures that can be tested now; broader scientific truth
and banking relevance require the substantive evidence that the workflow is
designed to help a reviewer inspect.

\label{exec17:validity-development-end}
